\chapter{Kết quả và đánh giá}
\label{chap:results}

Chương này trình bày phương pháp đánh giá hệ thống và các số liệu minh
họa cho năm RQ đã đặt ra ở Chương~\ref{chap:introduction}. Toàn bộ số
liệu trong chương này là \textbf{số liệu minh hoạ} dựa trên các benchmark
công bố gần đây~\cite{shufflebench_2024, fault_recovery_stream_2024} và
tài liệu kỹ thuật chính thức của Confluent~\cite{confluent_partition_count,
confluent_optimize_clients}. Mục tiêu là minh hoạ phương pháp phân tích
và kỳ vọng hợp lý khi triển khai thực tế trên cụm phát triển --- không
phải kết quả đo trực tiếp trên cluster đề xuất.

\section{Phương pháp benchmark}

\subsection{Thiết kế thực nghiệm}

Benchmark gồm tám kịch bản, mỗi kịch bản nhắm vào một hoặc nhiều RQ:

\begin{table}[H]
\centering
\caption{Mapping giữa kịch bản benchmark và RQ.}
\label{tab:benchmark-rq}
\begin{tabularx}{\textwidth}{l l l X}
\toprule
\textbf{Kịch bản} & \textbf{Biến độc lập} & \textbf{RQ} & \textbf{Mục đích} \\
\midrule
Smoke & rate 100--1k evt/s & --- & Xác nhận pipeline chạy đúng \\
Load ramp & rate 1k $\rightarrow$ 50k evt/s & RQ-A & Đo throughput \& ingestion lateness (cơ sở H1) \\
Partition scaling & partitions 3, 6, 12 & RQ-B & Ảnh hưởng partition count \\
Executor scaling & workers 1, 2, 4 & RQ-B & Ảnh hưởng số executor \\
Fault injection & kill worker tại T = 60s & RQ-C & Recovery time \& event loss \\
Veracity injection & 95/2/1/1/0.5/0.5\% & RQ-D & Dedup, late, bad record; kế toán record (H2) \\
Storage replay & 1, 5, 20 GB & RQ-E & Query latency, file count \\
Compaction & before / after OPTIMIZE & RQ-E & Tệp nhỏ $\rightarrow$ tệp lớn \\
\bottomrule
\end{tabularx}
\end{table}

\subsection{Metric thu thập}

Metric được chia thành mười ba nhóm:

\begin{table}[H]
\centering
\caption{Bộ metric thu thập theo nhóm.}
\label{tab:metrics}
\small
\begin{tabularx}{\textwidth}{l X l}
\toprule
\textbf{Nhóm} & \textbf{Metric} & \textbf{Nguồn} \\
\midrule
Ingestion &
  produced events/s, bytes/s, failed sends &
  Producer JMX, custom counter \\
Kafka &
  partition count, broker throughput, consumer lag &
  Kafka JMX, Kafka Exporter~\cite{kafka_docs_4_3} \\
Spark &
  numInputRows, processedRowsPerSecond, batchDuration &
  Spark query progress API~\cite{spark_docs_4_2} \\
Ingestion lag &
  Kafka consumer lag (LEO $-$ currentOffset) theo partition &
  Kafka Exporter \\
Ingestion lateness &
  p50/p95/p99 của $t^{s} - t^{e}$
  (\texttt{ingest\_time\_spark} $-$ \texttt{event\_time}) &
  Tính tại Bronze job từ hai timestamp có sẵn \\
Publication latency &
  p50/p95/p99 của $t^{o} - \mathrm{end}(W_m)$ trên các window đã chốt &
  Tính tại Gold job (cần thêm timestamp công bố; xem Mục~\ref{subsec:gold-schema}) \\
State &
  state rows, state memory, watermark lag &
  Spark metrics \\
Quality &
  p95 lateness, late-event ratio, duplicate ratio, quarantine ratio
  (Mục~\ref{subsec:quality-metrics}) &
  Bronze $\rightarrow$ Silver job \\
Reliability &
  lost events, duplicate events &
  So sánh số record theo \texttt{event\_id} \\
Recovery &
  time từ kill đến \texttt{processedRowsPerSecond} $\geq$ 90\% baseline &
  Spark query progress logs \\
Compute &
  CPU\%, RAM, GC time &
  JMX Exporter~\cite{prometheus_jmx_exporter} \\
HDFS &
  used space, read/write throughput &
  HDFS NameNode JMX~\cite{hdfs_docs_3_5} \\
Delta &
  file count, average file size &
  \texttt{DESCRIBE DETAIL} của Delta~\cite{delta_lake_docs} \\
\bottomrule
\end{tabularx}
\end{table}

\subsection{Công thức tính}
\label{subsec:formulas}

Báo cáo dùng một taxonomy nhất quán với Mục~\ref{subsec:event-tuple}: ba đại
lượng đo ``độ trễ'' khác nhau, không gộp chung thành một con số.

\textbf{Ingestion lateness} (trễ vào, mức event): theo đúng định nghĩa
Ct.~\ref{eq:observed-lateness},
\begin{equation}
  L_i = \max\bigl(0,\ t^{s}_i - t^{e}_i\bigr),
  \label{eq:ingestion-lateness}
\end{equation}
tính từ \texttt{ingest\_time\_spark} và \texttt{event\_time} của Bronze.
Đây là đại lượng đo được trực tiếp từ hai timestamp của từng event; nó
\emph{không} đồng nhất với lateness so với watermark: một event được coi là
``đến sau watermark'' khi tại thời điểm đến, watermark $W_Q$ của query đã
vượt qua event time của nó (Mục~\ref{subsec:watermark}) --- vị trí này phụ
thuộc cả tiến độ đọc $M_Q$ của query, nên $L_i$ vượt watermark bound chỉ là
chỉ báo nguy cơ, không tự động là post-watermark; và một event post-watermark
cũng không đồng nghĩa bị xoá (Mục~\ref{subsec:watermark}).

\textbf{Ingestion lag} (trễ vào, mức luồng):
\begin{equation}
  \mathrm{Lag}_{\text{group, topic, partition}} = \mathrm{LEO} - \mathrm{currentOffset}
  \label{eq:consumer-lag}
\end{equation}
trong đó LEO (Log End Offset) là vị trí cuối cùng của log Kafka, còn
\texttt{currentOffset} là vị trí consumer đang đọc tới. Lag đo độ trễ giữa
thời điểm bản ghi \emph{có trong Kafka} và thời điểm Spark đọc nó; khi input
rate liên tục vượt processing capacity, lag tăng tuyến tính thay vì tự tiêu
tan~\cite{confluent_optimize_clients}.

\textbf{Publication latency} (trễ ra, mức window): thời điểm một aggregate
xuất hiện ở Gold so với thời điểm cửa sổ đó kết thúc theo event time:
\begin{equation}
  \mathrm{PubLatency}_m = t^{o}_m - \mathrm{end}(W_m),
  \label{eq:publication-latency}
\end{equation}
với $t^{o}_m$ là thời điểm dòng aggregate của $W_m$ được ghi vào Gold. Vì Gold
job là window aggregate, publication latency tồn tại ở mức window, không phải
mức event (hệ quả nêu trong Mục~\ref{subsec:gold-schema}).

Với mỗi phân phối, báo cáo dùng p50, p95 và p99 để đặc trưng
~\cite{iot_stream_latency_2021}. Các bảng RQ-A dưới đây trình bày số minh hoạ
theo metric ingestion lateness (Ct.~\ref{eq:ingestion-lateness}) --- đại lượng
đo được trực tiếp từ hai trường timestamp sẵn có --- và gọi tên rõ để tránh
nhầm với publication latency.

Throughput tức thời:
\begin{equation}
  \mathrm{Throughput}(t) = \frac{N_{\text{processed}}(t)}{\Delta t}
  \label{eq:throughput}
\end{equation}
trong đó $N_{\text{processed}}(t)$ là số record được job xử lý
trong cửa sổ trượt $\Delta t$ (ví dụ 10 giây).

\subsection{Metric chất lượng (rolling)}
\label{subsec:quality-metrics}

Nhóm metric Quality trong Bảng~\ref{tab:metrics} cài đặt bốn trong sáu
đại lượng rolling của Mục~\ref{subsec:rolling-metrics}:
\begin{equation}
  \rho_{\mathrm{late}} = \frac{|\{e_i : L_i > \ell\}|}{N_{\mathrm{in}}},
  \qquad
  \rho_{\mathrm{dup}} = \frac{N_{\mathrm{dup,drop}}}{N_{\mathrm{in}}},
  \qquad
  \rho_{\mathrm{inv}} = \frac{N_{\mathrm{quar}}}{N_{\mathrm{in}}},
  \label{eq:quality-ratios}
\end{equation}
trong đó $\ell$ là ngưỡng late (baseline: $\ell = 10$ phút, trùng watermark
bound), $N_{\mathrm{dup,drop}}$ là số bản ghi bị dedup loại công khai, và
$N_{\mathrm{quar}}$ là số bản ghi vào \texttt{silver\_quarantine}. Lưu ý
ngữ nghĩa: $\ell$ đặt trên observed lateness $L_i$ là một ngưỡng vận hành
chọn được; nó không đồng nhất với vị từ post-watermark của query
(Mục~\ref{subsec:watermark}) --- một event có $L_i > \ell$ không tự động đến
sau watermark và ngược lại. Bốn đại
lượng này, cùng watermark lag (thuộc nhóm State), là đầu vào của policy
FAST/CAUTIOUS đề xuất (Mục~\ref{subsec:policy}) và của giả thuyết H2
(Mục~\ref{subsec:hypotheses}); đại lượng thứ sáu (policy transition count)
chỉ có ý nghĩa khi extension ở Mục~\ref{subsec:policy} được triển khai.

\section{Kết quả minh hoạ}

\textbf{Lưu ý:} tất cả số liệu dưới đây là \textbf{minh hoạ}, có đơn vị
và xu hướng phù hợp với~\cite{shufflebench_2024, fault_recovery_stream_2024,
confluent_partition_count, confluent_optimize_clients}, không phải số
đo trên cluster đề xuất. Khi triển khai thực tế, nhóm cần chạy script
\texttt{load\_test.sh} và ghi đè các giá trị trong bảng bằng dữ liệu thật.

\subsection{RQ-A: Throughput và ingestion lateness theo input rate}

Bảng~\ref{tab:rq-a} và Hình~\ref{fig:rq-a} cho thấy quan hệ giữa input
rate và ingestion lateness (Ct.~\ref{eq:ingestion-lateness}). Xu hướng minh
hoạ: dưới ngưỡng capacity ($\sim$20k evt/s), lateness tăng nhẹ theo rate;
trên ngưỡng, lateness tăng phi tuyến do micro-batch bị lấp đầy và consumer
lag tăng.

\begin{table}[H]
\centering
\caption{Throughput và ingestion lateness theo input rate (số liệu minh hoạ, đơn vị: event/s và ms).}
\label{tab:rq-a}
\begin{tabular}{rrrrr}
\toprule
\textbf{Input rate} & \textbf{Throughput} & \textbf{p50} & \textbf{p95} & \textbf{p99} \\
\midrule
1,000  & 1,000  & 850  & 1,420 & 2,310 \\
5,000  & 4,990  & 920  & 1,580 & 2,640 \\
10,000 & 9,950  & 1,050 & 1,890 & 3,120 \\
25,000 & 24,300 & 1,920 & 3,470 & 5,820 \\
50,000 & 38,600 & 6,840 & 14,200 & 23,500 \\
\bottomrule
\end{tabular}
\end{table}

\begin{figure}[H]
\centering
\begin{tikzpicture}
\begin{axis}[
    width=12cm, height=6cm,
    xlabel={Input rate (event/s)},
    ylabel={Ingestion lateness (ms)},
    xmode=log, ymode=log,
    grid=major, grid style={dashed, gray!30},
    legend pos=north west,
    legend style={font=\footnotesize},
    title={\small RQ-A: Ingestion lateness theo input rate (minh hoạ)}
]
\addplot[mark=*, thick, blue] coordinates {
    (1000, 850) (5000, 920) (10000, 1050) (25000, 1920) (50000, 6840)
};
\addlegendentry{p50}
\addplot[mark=square*, thick, orange] coordinates {
    (1000, 1420) (5000, 1580) (10000, 1890) (25000, 3470) (50000, 14200)
};
\addlegendentry{p95}
\addplot[mark=triangle*, thick, red] coordinates {
    (1000, 2310) (5000, 2640) (10000, 3120) (25000, 5820) (50000, 23500)
};
\addlegendentry{p99}
\end{axis}
\end{tikzpicture}
\caption{RQ-A: p99 ingestion lateness tăng phi tuyến khi input rate vượt
capacity của hệ thống. Số liệu minh hoạ
theo~\cite{shufflebench_2024, confluent_optimize_clients}.}
\label{fig:rq-a}
\end{figure}

\textbf{Nhận xét:} tại 50k evt/s, throughput thực tế (38,600 evt/s) thấp
hơn input, thể hiện qua lag Kafka tăng dần. Đây là dấu hiệu saturation
cần được phát hiện bằng alert \texttt{Spark processedRowsPerSecond <
inputRowsPerSecond}.

\subsection{RQ-B: Ảnh hưởng của partition count và worker count}

Bảng~\ref{tab:rq-b} minh hoạ ảnh hưởng của partition count đến throughput
tại cùng một input rate (25k evt/s). Tăng partition không phải lúc nào
cũng cải thiện tuyến tính~\cite{confluent_partition_count}.

\begin{table}[H]
\centering
\caption{Throughput và p99 ingestion lateness theo partition count (workload 25k evt/s, 2 worker, minh hoạ).}
\label{tab:rq-b}
\begin{tabular}{rrrr}
\toprule
\textbf{Partitions} & \textbf{Throughput} & \textbf{p99} & \textbf{CPU avg} \\
\midrule
3  & 18,200 & 8,950 & 78\% \\
6  & 23,400 & 4,120 & 71\% \\
12 & 24,800 & 2,980 & 64\% \\
\bottomrule
\end{tabular}
\end{table}

Hình~\ref{fig:rq-b} cho thấy tương tự cho số Spark worker (workload
25k evt/s, 6 partitions).

\begin{figure}[H]
\centering
\begin{tikzpicture}
\begin{axis}[
    width=11cm, height=6cm,
    xlabel={Số Spark worker},
    ylabel={Throughput (event/s)},
    ymin=10000, ymax=26000,
    grid=major, grid style={dashed, gray!30},
    title={\small RQ-B: Throughput theo số worker (25k evt/s, 6 partitions, minh hoạ)}
]
\addplot[mark=*, thick, blue] coordinates {
    (1, 11800) (2, 22400) (4, 24600)
};
\end{axis}
\end{tikzpicture}
\caption{RQ-B: Throughput tăng gần tuyến tính từ 1 đến 2 worker, sau đó
chững lại do overhead điều phối và bottleneck tại Kafka producer. Số liệu
minh hoạ theo~\cite{shufflebench_2024}.}
\label{fig:rq-b}
\end{figure}

\textbf{Nhận xét:} scaling từ 1 lên 2 worker cho throughput tăng gần gấp
đôi; từ 2 lên 4 worker tăng thêm chỉ $\sim$10\% vì bottleneck dịch chuyển
sang Kafka producer. Đây là một trong những trade-off điển hình của hệ
phân tán: scale từng thành phần phải đi kèm scale thành phần liên quan~\cite{stream_processing_evolution_2024}.

\subsection{RQ-C: Fault recovery khi kill worker}

Hình~\ref{fig:rq-c} mô tả diễn biến của consumer lag và throughput khi
một Spark worker bị kill tại $T = 60$\,s và ngừng xử lý cho đến khi được
khởi động lại tại $T = 180$\,s.

\begin{figure}[H]
\centering
\begin{tikzpicture}
\begin{axis}[
    width=12cm, height=6cm,
    xlabel={Thời gian $T$ (giây)},
    ylabel={Consumer lag (events)},
    grid=major, grid style={dashed, gray!30},
    title={\small RQ-C: Consumer lag khi kill worker tại T = 60s (minh hoạ)}
]
\addplot[thick, blue] coordinates {
    (0,0) (15,120) (30,160) (45,140) (60,180)
};
\addlegendentry{pre-failure (T $<$ 60)}
\addplot[thick, red, dashed] coordinates {
    (60,180) (90,7800) (120,15600) (150,20400) (180,24000)
};
\addlegendentry{outage (60 $\leq$ T $<$ 180)}
\addplot[thick, orange] coordinates {
    (180,24000) (195,26500) (210,21000) (225,12000) (240,5000)
};
\addlegendentry{restart \& drain (180 $\leq$ T $<$ 240)}
\addplot[thick, green!50!black] coordinates {
    (240,5000) (255,2200) (270,900) (285,350) (300,120)
};
\addlegendentry{post-recovery (T $\geq$ 240)}
\end{axis}
\end{tikzpicture}
\caption{RQ-C: Diễn biến consumer lag qua bốn giai đoạn phân tách theo thời
điểm: (i) hoạt động bình thường $T < 60$\,s; (ii) worker bị kill tại
$T = 60$\,s, pipeline ngừng xử lý và lag tích luỹ gần tuyến tính đến
$T = 180$\,s; (iii) worker được khởi động lại tại $T = 180$\,s, lag tạm tăng
nhẹ khi khôi phục checkpoint rồi bắt đầu giảm trong lúc backlog được xử lý;
(iv) từ $T \geq 240$\,s lag giảm nhanh và về mức trước failure ($\sim$120
events) tại $T \approx 300$\,s. Số liệu
minh hoạ theo~\cite{fault_recovery_stream_2024}.}
\label{fig:rq-c}
\end{figure}

\textbf{Metric recovery time:} theo định nghĩa ở Bảng~\ref{tab:metrics},
tính từ thời điểm kill ($T = 60$\,s trên Hình~\ref{fig:rq-c}) đến khi
\texttt{processedRowsPerSecond} đạt $\geq 90\%$ baseline. Trong minh hoạ,
worker chỉ được khởi động lại thủ công tại $T = 180$\,s, nên recovery time
$\approx 180$\,s, trong đó 120\,s là downtime do chính sách restart thủ công
(không phải chi phí của hệ thống) và phần còn lại là khôi phục checkpoint cùng
ấm lại xử lý --- nhất quán với kết quả ``Spark phục hồi ổn định nhờ checkpoint
nhưng có đánh đổi latency trong giai đoạn restart'' của
~\cite{fault_recovery_stream_2024}. Consumer lag cần thêm thời gian tiêu tan
backlog và mới về mức trước failure ($\sim$120 events) tại $T \approx
300$\,s, tức phục hồi hoàn toàn muộn hơn thời điểm throughput hồi phục.

\textbf{Lost / duplicate events:} nhờ idempotent producer (\texttt{enable.idempotence}
và \texttt{acks=all}) cùng Delta transaction log, số event bị mất được
kỳ vọng bằng 0. Số duplicate ở Silver được kiểm soát bởi
\texttt{dropDuplicates} trong Bronze job.

\subsection{RQ-D: Data veracity --- dedup, late, bad record}

Bảng~\ref{tab:rq-d} cho thấy hiệu quả của các cơ chế xử lý veracity
khi producer cố tình inject 5\% record không hoàn hảo. Bảng tuân thủ bất
biến kế toán I1 (Ct.~\ref{eq:accounting}): với mỗi loại, cột Bronze in bằng
tổng Silver out + Quarantine + Explicitly dropped; toàn bộ 100.000 record đầu
vào đều được giải trình, không có record ``mất tích''.

\begin{table}[H]
\centering
\caption{Kế toán record khi chạy với veracity injection profile (số liệu minh hoạ).
``Explicit drop'' là bản trùng bị loại bởi dedup theo \texttt{event\_id}
(pass 1 ở Bronze trong watermark bound; bản trùng đến muộn qua bound bị pass 2
của Silver batch loại); ``Quarantine'' là bản ghi vi phạm required-field hoặc
physical bounds được Silver job tách kèm lý do.}
\label{tab:rq-d}
\begin{tabular}{l r rrrr}
\toprule
\textbf{Loại record} & \textbf{Inject} & \textbf{Bronze in} & \textbf{Silver out} & \textbf{Quarantine} & \textbf{Explicit drop} \\
\midrule
Valid          & 95.0\% & 95,000 & 95,000 & 0    & 0 \\
Duplicate      & 2.0\%  & 2,000  & 0      & 0    & 2,000 \\
Missing field  & 1.0\%  & 1,000  & 0      & 1,000 & 0 \\
Out of range   & 1.0\%  & 1,000  & 0      & 1,000 & 0 \\
Late 30s       & 0.5\%  & 500    & 500    & 0    & 0 \\
Late 5min      & 0.5\%  & 500    & 500    & 0    & 0 \\
\midrule
\textbf{Total} & 100\%  & 100,000 & 96,000 & 2,000 & 2,000 \\
\bottomrule
\end{tabular}
\end{table}

\textbf{Nhận xét.} Hai nghìn bản trùng bị loại công khai bởi dedup theo
\texttt{event\_id}: pass 1 ở Bronze loại các bản trùng đến trong watermark
bound, còn pass 2 (batch, trong Silver) là lớp bảo vệ thứ hai cho bản trùng
đến muộn sau khi state của pass 1 đã được watermark giải phóng; mỗi bản trùng
bị loại đúng một lần, nên bất biến I1 không bị vi phạm. Missing field và
out-of-range được Silver tách sang \texttt{silver\_quarantine} kèm
\texttt{quarantine\_reason} chứ không bị xoá. Hai loại late (30s và 5min) có
$L_i$ nhỏ hơn watermark bound baseline 10 phút; với luồng đến liên tục,
watermark tại thời điểm chúng đến thường chưa vượt qua event time của chúng,
nên chúng được ingest bình thường vào Bronze/Silver và, nếu đến trước khi cửa
sổ tương ứng được chốt, vẫn góp vào aggregate của cửa sổ --- watermark không
xoá event muộn. Điều kiện này mang tính xác suất theo tiến độ đọc $M_Q$ của
query chứ không phải kết luận từ một event riêng lẻ: như phân tích ở
Mục~\ref{subsec:watermark}, vị trí so với watermark không suy ra trực tiếp
từ $L_i$. Event đến sau watermark mà cửa sổ Gold của nó đã được công bố thì
trong baseline \emph{không} tự động góp vào aggregate đã chốt: Gold job chạy
ở append output mode có watermark, engine lọc record cũ hơn watermark khỏi
phép tổng hợp và không phát hành thêm dòng cho cửa sổ đã chốt; event đó cũng
không tự động vào quarantine. Baseline hiện tại không có cơ chế sửa kết quả;
việc phát hành lại giá trị cửa sổ (CORRECTED) là mục tiêu của vòng đời
PROVISIONAL/FINAL/CORRECTED kèm reconciliation đề xuất
(Mục~\ref{subsec:lifecycle}). Con số cụ thể là minh hoạ; kiểm chứng định
lượng thuộc giả thuyết H2 (Mục~\ref{subsec:hypotheses}).

\subsection{RQ-E: Delta partitioning và compaction}

Bảng~\ref{tab:rq-e} so sánh query latency và số tệp Delta trước và sau
\texttt{OPTIMIZE} với dữ liệu replay 5 GB.

\begin{table}[H]
\centering
\caption{Ảnh hưởng của OPTIMIZE lên query latency và số tệp (5 GB replay, minh hoạ).
Q1/Q3/Q5 là ba truy vấn analytical mẫu được định nghĩa ngay dưới bảng.}
\label{tab:rq-e}
\begin{tabular}{lrr rrr}
\toprule
\textbf{Trạng thái} & \textbf{Tệp} & \textbf{Avg file size (MB)} & \textbf{Q1} & \textbf{Q3} & \textbf{Q5} \\
\midrule
Trước OPTIMIZE  & 8,420 & 0.59  & 4.2s & 12.8s & 18.5s \\
Sau OPTIMIZE    & 64    & 78.0  & 0.6s & 1.4s  & 2.1s  \\
\bottomrule
\end{tabular}
\end{table}

trong đó Q1, Q3, Q5 là ba truy vấn analytical mẫu (count theo location,
average value theo zone, top-10 sensor có variance cao). Việc gộp hàng
ngàn tệp nhỏ ($\sim$0.59 MB) thành hàng chục tệp lớn ($\sim$78 MB) cải
thiện query latency khoảng 7--9 lần trong minh hoạ này~\cite{delta_lake_docs}.

\subsection{Khả năng đo được trên schema hiện tại}
\label{subsec:gold-schema}

Hai metric trong Bảng~\ref{tab:metrics} cần lưu ý về tính khả thi trên schema
Gold hiện tại:

\begin{itemize}
  \item \textbf{Publication latency} (Ct.~\ref{eq:publication-latency}) đòi hỏi
        mỗi dòng aggregate mang thời điểm được ghi vào Gold. Bản cài đặt hiện
        tại không ghi trường này, nên metric chưa đo được trực tiếp; muốn đo,
        cần bổ sung một cột timestamp công bố (ví dụ \texttt{current\_timestamp()}
        trong Gold job) và chạy lại --- thay đổi nhỏ ở tầng Gold nhưng là thay
        đổi thật, không phải cấu hình.
  \item \textbf{Anomaly count} xuất hiện trong một số mô tả bảng tổng hợp
        trước đây nhưng \emph{không} nằm trong schema Gold hiện tại
        (avg/min/max/count theo cửa sổ). Báo cáo không định nghĩa
        per-event output latency hay anomaly count cho baseline; chúng chỉ có
        thể bàn như một phần của extension.
\end{itemize}

\section{Phân tích và thảo luận}

\subsection{Mapping giả thuyết với kịch bản đo}
\label{subsec:hypothesis-mapping}

Bốn giả thuyết ở Mục~\ref{subsec:hypotheses} được gắn với kịch bản thực
nghiệm cụ thể để khi có cụm thật, mỗi giả thuyết có đường đo riêng.
Không giả thuyết nào được xem là đã xác nhận bởi số liệu minh hoạ của chương
này.

\begin{table}[H]
\centering
\caption{Mapping giữa giả thuyết và kịch bản/metric cần đo.}
\label{tab:hypothesis-mapping}
\begin{tabularx}{\textwidth}{l l X}
\toprule
\textbf{Giả thuyết} & \textbf{Kịch bản} & \textbf{Metric quyết định} \\
\midrule
H1 (allowed lateness) & Load ramp + chạy lặp với $\theta \in \{1, 5, 10, 30\}$ phút &
  retention của late event, state rows/state memory, publication delay \\
H2 (kế toán chất lượng) & Veracity injection &
  $\rho_{\mathrm{inv}}, \rho_{\mathrm{dup}}, \rho_{\mathrm{late}}$ so với tỷ lệ
  inject, kèm sai số lấy mẫu \\
H3 (reconciliation) & Storage replay + batch correction job (extension) &
  sai lệch aggregate trước/sau chỉnh sửa trên các window đủ cũ \\
H4 (adaptive policy) & Load ramp có shift chất lượng kéo dài &
  ablation fixed vs adaptive, adaptive có/không hysteresis (Ct.~\ref{eq:hysteresis}) \\
\bottomrule
\end{tabularx}
\end{table}

\subsection{Trade-off chính}

Hệ thống đề xuất thể hiện ba trade-off điển hình của Big Data Systems:

\begin{enumerate}
  \item \textbf{Throughput vs ingestion lateness:} trigger interval nhỏ (1--2s)
        giảm ingestion lateness nhưng tăng overhead mỗi micro-batch; trigger
        lớn (30--60s) ngược lại. Cấu hình hiện tại dùng 5s cho Bronze và 10s cho Gold ---
        cân bằng hợp lý cho workload 25k evt/s~\cite{shufflebench_2024};
  \item \textbf{Partition count vs file overhead:} partition nhiều tăng
        parallelism nhưng cũng tăng số connection Kafka, số tệp Delta
        nhỏ, áp lực lên NameNode. Kết quả cho thấy 6 partitions phù
        hợp với 2 worker trong phạm vi phát triển~\cite{confluent_partition_count};
  \item \textbf{Watermark bound vs state size:} watermark lớn giữ được
        nhiều late event nhưng tăng state và làm chậm publication
        (nội dung giả thuyết H1); watermark nhỏ giải phóng
        state nhanh nhưng để lọt nhiều late event khỏi dedup. 10 phút là giá
        trị baseline minh hoạ cân bằng giữa hai yếu
        tố~\cite{iot_stream_latency_2021}.
\end{enumerate}

\subsection{Giới hạn}

Số liệu trong chương này có ba giới hạn cần ghi nhận:

\begin{itemize}
  \item \textbf{Chưa đo trên cluster thật}: số liệu minh hoạ được lấy
        theo xu hướng từ benchmark công bố, không phải đo trực tiếp.
        Sai lệch $\pm 30\%$ so với thực tế là có thể;
  \item \textbf{Cluster nhỏ}: cụm Docker Compose dùng để phát triển có
        1 Kafka broker, 2 worker, 2 DataNode --- không đại diện cho môi
        trường production;
  \item \textbf{Không có HA}: cả NameNode lẫn Kafka broker đều ở chế độ
        single-node. Khi chuyển sang production cần HDFS HA với QJM và
        Kafka replication $\geq 3$.
\end{itemize}

\subsection{Khuyến nghị cho triển khai thực tế}

\begin{itemize}
  \item Chạy benchmark baseline trước khi thay đổi bất kỳ cấu hình nào,
        lưu kết quả vào \texttt{10\_Results/} để so sánh;
  \item Sử dụng baseline đã đo được làm threshold cho Grafana alert,
        tránh đặt ngưỡng "copy từ Internet"~\cite{grafana_alerting};
  \item Ưu tiên tối ưu Kafka producer (batching, compression) trước
        khi tăng số executor~\cite{confluent_optimize_clients};
  \item Đặt lịch \texttt{OPTIMIZE} hàng ngày cho các bảng streaming để
        giải quyết small-file problem~\cite{delta_lake_docs}.
\end{itemize}

\section{Tóm tắt chương}

Chương này đã trình bày:
\begin{itemize}
  \item Phương pháp benchmark gồm tám kịch bản mapping tới năm RQ và bốn giả
        thuyết H1--H4;
  \item Bộ metric mười ba nhóm với taxonomy nhất quán (ingestion lateness,
        ingestion lag, publication latency), công thức throughput và metric
        chất lượng rolling;
  \item Số liệu minh hoạ cho từng RQ cùng phân tích trade-off, với bảng
        kế toán record bảo toàn cho RQ-D;
  \item Khuyến nghị khi triển khai thực tế.
\end{itemize}

Toàn bộ kết luận định lượng ở đây vẫn là minh hoạ; việc xác nhận hay bác bỏ
H1--H4 chỉ có thể thực hiện khi chạy benchmark thật trên cụm đề xuất.
Chương tiếp theo tổng kết báo cáo và đề xuất hướng phát triển.